% Conclusion: direct answer to the project's question

\section{Discussion: how hard is the problem?}\label{sec:conclusion}

The project asked how hard cardinality-constrained portfolio selection is empirically.
On real S\&P~500 data the answer splits cleanly along the in-sample / out-of-sample axis,
and the two halves of the answer point in opposite directions.

\paragraph{In-sample, the problem is computationally easy.} A deterministic spectral seed
plus a 1-swap polish (\textsc{Slh}) matches the brute-force oracle on every cell of the
$(n,k)$ grid we can verify by enumeration, at a cost of roughly $0.1$\,ms per call. Greedy
forward selection stays within $0.05$ to $0.12$\,bp of the oracle; simulated annealing
matches the oracle on small cells and degrades gently on larger ones; \textsc{Slh}\
combines the strengths of both deterministically. The interesting combinatorial behaviour
the OGP literature warns about, where convex relaxations stall in a forbidden-overlap
plateau, is not the empirical bottleneck here. The reason is structural: Phase~1
documented a strong rank-one market factor in $\hat\Sigma$ together with a near-isotropic
residual, and Phase~4a showed that this is exactly the kind of factor-plus-residual
covariance under which top-$k$ thresholding of a ridge-whitened $\hat\mu$ provably
recovers the planted optimum (\cref{prop:slh-recovery}). The empirical SNR sits in this
recovery regime, not the OGP-active regime.

\paragraph{Out-of-sample, the problem is statistically hard.} This is the main practical
finding of the project. The brute-force oracle achieves in-sample Sharpe between $0.91$
and $1.16$ across the $(n,k)$ cells we test; on the $20\%$ held-out block of the
same panel its out-of-sample Sharpe drops to $0.46$ to $0.71$, a relative collapse of
$30$ to $50\%$. Greedy, simulated annealing, and \textsc{Slh}\ inherit identical
collapses, because they all recover the same support as the oracle. Finding a better
in-sample optimum simply does not buy out-of-sample performance.

This is not a failure of the optimisation pipeline. It is a sample-size limitation in
$\hat\mu$ and $\hat\Sigma$. Phase~1 placed the residual variance proxy at roughly $5\times$
the Gaussian benchmark (from $\psi_2/\nu_{\mathrm{Gauss}}\approx 2.2$); for $T=2{,}009$
training days the per-coordinate signal-to-noise ratio of $\hat\mu_i$ is bounded by a
constant of order one, and Phase~3a's threshold $a_\star(\rho)$ is essentially a count of
the standard deviations separating the planted signal from the bulk. With realistic
financial data this margin is small, which means the in-sample top-$k$ is genuinely
informative about \emph{some} signal but most of that signal is sample noise
specific to the training block. The 30 to 50\,\% Sharpe collapse is the discrete-portfolio
analogue of the bias-variance trade-off familiar from high-dimensional regression: the
in-sample optimum is over-fit to the realisation of $\hat\mu,\hat\Sigma$.

\paragraph{The role of the OGP theory in this story.} The Phase~3 theorems and the
algorithmic hardness corollary are not directly responsible for the in-sample performance
of \textsc{Slh}\ on real data: as discussed in \S\ref{ssec:slh-regime}, the empirical
SNR sits above the recovery threshold, where the OGP-active regime is not relevant. The
theory does two other useful things. First, it quantifies the boundary above which any
algorithm could conceivably recover the planted optimum (\cref{thm:detect}). Second, it
identifies an explicit forbidden-overlap interval (\cref{thm:forbid}) for the linear
score and a conjectural classical OGP for the full quadratic objective, supported by
the mixing-time blow-up of constrained Metropolis between $(n,k)=(20,5)$ and
$(40,10)$ in \S\ref{ssec:p2b}. If we ever moved to a regime with weaker per-asset
signal, this is the territory the algorithm would have to enter.

\paragraph{Where the project would go next.} Three directions, in decreasing order of
practical importance.
\begin{itemize}\itemsep=2pt
\item \emph{Statistical regularisation of $\hat\mu,\hat\Sigma$.} A Ledoit-Wolf shrinkage
estimator \cite{ledoit-wolf}, or a Bayesian factor-model prior, would attack the
out-of-sample collapse directly. Whether shrinkage moves the empirical SNR towards
$a_\star(\rho)$ (turning the problem from over-fit but combinatorially easy into honest
but combinatorially harder) is a question the regime framework of \S\ref{ssec:slh-regime}
makes concrete.
\item \emph{A second-moment proof of the classical OGP for $J$.} The empirical
mixing-time blow-up in \S\ref{ssec:p2b} is suggestive but not definitive. A
Sherrington-Kirkpatrick-style second-moment computation on the spike-plus-bulk model of
\eqref{eq:factor-noise} would convert the conjecture into a theorem. The relevant
machinery is the Gaussian-comparison/Lindeberg toolkit cited in \S\ref{sec:ogp}.
\item \emph{Risk-aversion sensitivity and richer experiments.} We fixed $\gamma=10$
throughout. A sensitivity analysis at $\gamma\in\{1,10,100\}$ would test whether the
in-sample/out-of-sample story is robust to this choice. The Metropolis OGP evidence in
\S\ref{ssec:p2b} could be strengthened with multiple computational budgets, full trace
plots, and a comparison of cold-start versus planted-start hitting times; we report
basic hitting-time statistics here and leave the more thorough characterisation as
follow-up work.
\end{itemize}

\paragraph{Reproducibility.}
All numbers and figures come from the seven notebooks
\texttt{extract\_data.ipynb}, \texttt{01\_diagnostics.ipynb},
\texttt{02\_oracle\_baselines.ipynb}, \texttt{03\_mcmc\_sampler.ipynb},
\texttt{04a\_ogp\_theory.ipynb}, \texttt{04\_ogp\_experiment.ipynb}, and
\texttt{05\_hybrid\_algorithm.ipynb}, with deterministic seeds throughout. Total
wall-clock to regenerate every figure on a single laptop CPU is under five minutes.
